
\documentclass[ukrainian]{ucaphyslab}
\labtitle{Маятник Максвелла}
\labsubtitle{Лабораторна робота № (вкажіть номер)}

\student{Прізвище Ім'я}

\teacher{Володимир Коломієць} % якщо рядок не потрібен — видаліть текст і залиште \teacher{}

\labdate{РРРР-ММ-ДД}

\usepackage{float}
\usepackage{pgfplots}
\pgfplotsset{compat=1.18}
\definecolor{series1}{HTML}{2A78D6}
\definecolor{series2}{HTML}{EB6834}

\begin{document}
\maketitlepage

\section{Мета роботи}
Мета лабораторної роботи: ознайомитися з принципом дії маятника Максвелла. Дослідити динаміку плоского руху твердого тіла на прикладі
маятника Максвела. Обчислити та експериментально визначити момент інерції маятника.


\section{Прилади та обладнання}
\begin{itemize}
  \item Маятник Максвелла;
  \item
\end{itemize}

\section{Теоретичні відомості}

\textbf{Поняття та величини, що використовуються в роботі:}
\begin{itemize}
  \item \textbf{Імпульс частинки:} $\vec p = m\vec\upsilon$.
  \item \textbf{Момент імпульсу частинки:} $\vec L = [\vec r\times\vec p] = m[\vec r\times\vec\upsilon]$.
  \item \textbf{Момент сили:} $\vec M = [\vec r\times\vec F]$.
  \item \textbf{Кінетична енергія частинки:} $T = \dfrac{m\upsilon^2}{2}$.
  \item \textbf{Система центру мас (Ц-система)} — система відліку, початок якої знаходиться у центрі мас системи частинок; у ній система частинок як ціле перебуває в спокої.
  \item \textbf{Плоский рух твердого тіла} — рух, при якому всі точки тіла рухаються в паралельних площинах.
  % TODO: інерція, момент інерції (місце під тензори, формули 8.17–8.24 методички)
\end{itemize}

\textbf{Зміна імпульсу та моменту імпульсу частинки:}
\begin{equation}
  \frac{d\vec p}{dt} = \vec F, \tag{8.1}
\end{equation}
\begin{equation}
  \frac{d\vec L}{dt} = \vec M \quad \text{(рівняння моментів)}, \tag{8.3}
\end{equation}
де $\vec F$ — сума прикладених сил, $\vec M$ — сумарний момент сил.

\textbf{Для системи частинок} змінюють імпульс і момент імпульсу лише зовнішні сили:
\begin{equation}
  \frac{d\vec P}{dt} = \vec F^{\,\text{зовн}}, \qquad \frac{d\vec L}{dt} = \vec M^{\,\text{зовн}}. \tag{8.7, 8.12}
\end{equation}

\begin{keyfacts}
\textbf{Умови збереження:} імпульс зберігається, якщо сума (зовнішніх) сил дорівнює нулю; момент імпульсу зберігається, якщо сумарний момент (зовнішніх) сил дорівнює нулю.
\end{keyfacts}

\textbf{Зв'язок лінійних і кутових величин при обертальному русі:}
\begin{align}
  d\vec r &= [d\vec\varphi\times\vec r], \tag{8.14}\\
  \vec\upsilon &= [\vec\omega\times\vec r], \tag{8.15}\\
  \vec a &= [\vec\beta\times\vec r] + [\vec\omega\times[\vec\omega\times\vec r]], \tag{8.16}
\end{align}
де перший доданок у (8.16) — тангенціальна складова прискорення, другий — нормальна.

\section{Хід роботи та результати вимірювань}

У роботі вимірюють час падіння $\tau_i$ маятника з кожним кільцем і, знаючи радіус намотки
нитки $r_c$, за формулою (8.30) експериментально визначають момент інерції $I$.
Результат порівнюють з теоретичним значенням для маятника, що складається зі стержня А,
диска Д з отвором і кільця К (формули (8.33)–(8.36)).

\begin{table}[h]
  \centering
  \begin{tabular}{@{}llc@{}}
    \toprule
    Величина & Позначення & Значення \\
    \midrule
    Прискорення вільного падіння & $g$ & \SI{9.81}{\meter\per\second\squared} \\
    Висота падіння & $h$ & \SI{0.897}{\meter} \\
    Радіус стержня & $r_c$ & \SI{4.0}{\milli\meter} \\
    Радіус диска & $r_\text{д}$ & \SI{40.0}{\milli\meter} \\
    Внутрішній радіус кільця & $r_{\text{к}1}$ & \SI{40.0}{\milli\meter} \\
    Зовнішній радіус кільця & $r_{\text{к}2}$ & \SI{58.0}{\milli\meter} \\
    Маса стержня & $m_c$ & \SI{186.5}{\gram} \\
    Маса диска & $m_\text{д}$ & \SI{85.0}{\gram} \\
    Маса кільця 1 & $m_\text{к}$ & \SI{297.0}{\gram} \\
    Маса кільця 2 & $m_\text{к}$ & \SI{489.0}{\gram} \\
    \bottomrule
  \end{tabular}
  \caption{Константи та параметри установки}
\end{table}

\begin{table}[h]
  \centering
  \begin{tabular}{@{}c ccc ccc@{}}
    \toprule
     & \multicolumn{3}{c}{Кільце 1} & \multicolumn{3}{c}{Кільце 2} \\
    \cmidrule(lr){2-4} \cmidrule(lr){5-7}
    № & $\tau_1$, с & $\tau_2$, с & $\tau_\text{сер}$, с & $\tau_1$, с & $\tau_2$, с & $\tau_\text{сер}$, с \\
    \midrule
    1  & 3.36 & 3.84 & 3.600 & 4.18 & 4.07 & 4.125 \\
    2  & 3.66 & 3.76 & 3.710 & 3.65 & 3.76 & 3.705 \\
    3  & 3.33 & 3.49 & 3.410 & 3.45 & 3.54 & 3.495 \\
    4  & 3.51 & 3.52 & 3.515 & 3.50 & 3.64 & 3.570 \\
    5  & 3.46 & 3.39 & 3.425 & 3.35 & 3.72 & 3.535 \\
    6  & 3.61 & 3.56 & 3.585 & 3.36 & 3.59 & 3.475 \\
    7  & 3.45 & 3.32 & 3.385 & 3.50 & 3.69 & 3.595 \\
    8  & 3.26 & 3.59 & 3.425 & 3.43 & 3.64 & 3.535 \\
    9  & 3.15 & 3.34 & 3.245 & 3.53 & 3.57 & 3.550 \\
    10 & 3.21 & 3.31 & 3.260 & 3.61 & 3.71 & 3.660 \\
    \midrule
    Сер. & 3.400 & 3.512 & 3.456 & 3.556 & 3.693 & 3.625 \\
    \bottomrule
  \end{tabular}
  \caption{Час падіння маятника}
\end{table}

\section{Обробка результатів}
% Розрахунки шуканих величин та оцінка похибок вимірювань.
Експериментальний момент інерції:
\begin{equation}
  I_\text{експ} = m r_c^2\left(\frac{g\tau^2}{2h} - 1\right), \qquad m = m_c + m_\text{д} + m_\text{к}. \tag{8.30}
\end{equation}
Теоретичний момент інерції:
\begin{gather}
  I_\text{к} = \tfrac12 m_\text{к}\left(r_{\text{к}1}^2 + r_{\text{к}2}^2\right), \quad
  I_\text{д} = \tfrac12 m_\text{д}\left(r_c^2 + r_\text{д}^2\right), \quad
  I_c = \tfrac12 m_c r_c^2, \tag{8.33--8.35}\\
  I_\text{теор} = I_\text{к} + I_\text{д} + I_c. \tag{8.36}
\end{gather}
Розбіжність: $\Delta I = I_\text{експ} - I_\text{теор}$, \ $\varepsilon = \Delta I / I_\text{експ}\cdot 100\%$.

\begin{table}[H]
  \centering
  \begin{tabular}{@{}lcc@{}}
    \toprule
    Величина & Кільце 1 & Кільце 2 \\
    \midrule
    $\tau_\text{сер}$, с & \num{3.456} & \num{3.625} \\
    $m$, кг & \num{0.5685} & \num{0.7605} \\
    $I_\text{к}$, кг$\cdot$м$^2$ & \num{7.372e-4} & \num{1.214e-3} \\
    $I_\text{д}$, кг$\cdot$м$^2$ & \num{6.868e-5} & \num{6.868e-5} \\
    $I_c$, кг$\cdot$м$^2$ & \num{1.492e-6} & \num{1.492e-6} \\
    $I_\text{теор}$, кг$\cdot$м$^2$ & \num{8.073e-4} & \num{1.284e-3} \\
    $I_\text{експ}$, кг$\cdot$м$^2$ & \num{5.850e-4} & \num{8.619e-4} \\
    $\Delta I$, кг$\cdot$м$^2$ & \num{-2.223e-4} & \num{-4.219e-4} \\
    $\varepsilon$, \% & \num{-38.0} & \num{-49.0} \\
    \bottomrule
  \end{tabular}
  \caption{Результати розрахунків}
\end{table}

\begin{figure}[H]
  \centering
  \begin{tikzpicture}
    \begin{axis}[
      width=0.9\linewidth, height=6cm,
      xlabel={№ спроби}, ylabel={$\tau_\text{сер}$, с},
      xmin=0.5, xmax=10.5, xtick={1,...,10},
      ymin=3.1, ymax=4.2,
      ymajorgrids, grid style={gray!25},
      axis lines*=left,
      legend style={draw=none, at={(0.5,1.02)}, anchor=south, legend columns=2, font=\small},
    ]
      \addplot[color=series1, thick, mark=*, mark size=2pt] coordinates {
        (1,3.600) (2,3.710) (3,3.410) (4,3.515) (5,3.425)
        (6,3.585) (7,3.385) (8,3.425) (9,3.245) (10,3.260)};
      \addlegendentry{Кільце 1}
      \addplot[color=series2, thick, mark=square*, mark size=2pt] coordinates {
        (1,4.125) (2,3.705) (3,3.495) (4,3.570) (5,3.535)
        (6,3.475) (7,3.595) (8,3.535) (9,3.550) (10,3.660)};
      \addlegendentry{Кільце 2}
      \addplot[color=series1, dashed, forget plot] coordinates {(0.5,3.456) (10.5,3.456)};
      \addplot[color=series2, dashed, forget plot] coordinates {(0.5,3.625) (10.5,3.625)};
    \end{axis}
  \end{tikzpicture}
  \caption{Час падіння маятника в кожній спробі (штрихова лінія — середнє)}
\end{figure}

\begin{figure}[H]
  \centering
  \begin{tikzpicture}
    \begin{axis}[
      width=0.75\linewidth, height=6cm,
      ybar=2pt, bar width=18pt,
      symbolic x coords={Кільце 1, Кільце 2}, xtick=data,
      enlarge x limits=0.5,
      ylabel={$I$, $10^{-3}$ кг$\cdot$м$^2$},
      ymin=0, ymax=1.5,
      ymajorgrids, grid style={gray!25},
      axis lines*=left,
      nodes near coords, nodes near coords style={font=\small, color=black},
      every node near coord/.append style={/pgf/number format/.cd, fixed, precision=3, use comma},
      legend style={draw=none, at={(0.5,1.02)}, anchor=south, legend columns=2, font=\small},
    ]
      \addplot[fill=series1, draw=none] coordinates {(Кільце 1,0.585) (Кільце 2,0.862)};
      \addlegendentry{$I_\text{експ}$}
      \addplot[fill=series2, draw=none] coordinates {(Кільце 1,0.807) (Кільце 2,1.284)};
      \addlegendentry{$I_\text{теор}$}
    \end{axis}
  \end{tikzpicture}
  \caption{Експериментальний і теоретичний моменти інерції}
\end{figure}

\section{Висновки}
Сформулюйте висновок: чи узгоджується отриманий результат із теоретичним
(табличним) значенням у межах похибки, та які фактори могли вплинути на
точність вимірювань.

% Розкоментуйте, якщо цитуєте джерела через \cite{} (записи — у references.bib):
% \ucuprintbibliography

\end{document}
